% titel: Blatt 8, A5: Lokale Extrema und Sattelpunkte
% skript: Kap. 4.2--4.3 (Extrema, Sattelpunkte, Wendepunkte)

% Teil I: Werkzeuge, kein Lösungsweg.

Zwei Sätze über lokale Extrema, Sattel- und Wendepunkte beweisen.

\begin{werkzeugkasten}[Definitionen]
\begin{itemize}
  \item \textbf{Lokales Extremum} von $f$ bei $x_0$: $f$ nimmt dort einen
        lokalen Maximums-/Minimumswert an (Wert in einer Umgebung nicht
        größer/kleiner).
  \item \textbf{Isoliertes} Extremum: es gibt eine Umgebung, in der kein
        weiterer Extremalpunkt liegt (das Extremum ist ``getrennt'').
  \item \textbf{Wendepunkt}: Krümmungswechsel bei $x_0$; $f'$ hat dort ein
        lokales Extremum (Krümungsfunktion $f''$ wechselt das Vorzeichen).
  \item \textbf{Sattelpunkt}: $f'(x_0)=0$, aber kein Extremum; die
        Seitenstücke verlaufen monoton ($f$ ``biegt um, ohne zu wenden'' in
        vollem Sinne).
\end{itemize}
\textit{Skript:} Kap. 4.2--4.3.
\end{werkzeugkasten}

\begin{werkzeugkasten}[Kriterium über höhere Ableitungen]
Falls $f^{(n)}(x_0)=0$ für alle $n<m$ und $f^{(m)}(x_0)\ne0$, so entscheidet
die Parität von $m$:
\begin{itemize}
  \item $m$ gerade: $x_0$ ist isoliertes lokales Extremum (die
        Taylor-Entwicklung beginnt mit $\frac{f^{(m)}(x_0)}{m!}(x-x_0)^m$
        mit gerader Potenz, Vorzeichenwechsel gilt nicht).
  \item $m$ ungerade: $x_0$ ist isolierter Sattelpunkt (ungerade Potenz
        wechselt das Vorzeichen, aber nicht das Vorzeichen von $f'$ symmetrisch).
\end{itemize}

\textit{Skript:} Kap. 4.3.
\end{werkzeugkasten}

\begin{achtung}[Darauf achten]
\begin{itemize}
  \item Einen Wendepunkt über $f'$: $f'$ muss ein \emph{isoliertes} lokales
        Extremum haben; das ist die Formulierung der Krümmungsänderung.
  \item ``Isoliert'' ist essenziell: konstante Ableitung auf einem
        Intervall macht keinen Wendepunkt (kein Krümmungswechsel).
  \item Beim Beweis per Induktion die Fallunterscheidung gerade/ungerade
        sauber führen und den Induktionsanfang (niedrigste $m$) prüfen.
  \item Die Taylor-Vorzeichenargumentation statt bloßer
        ``$f'=0$ heißt Sattel/Estrum'' anwenden.
\end{itemize}
\end{achtung}
